\section{Experimental Setup}
\label{sec:baselines}

All numbers are measured on one NVIDIA TITAN RTX with
\texttt{GSAI-ML/LLaDA-8B-Instruct} in fp16. Prompts are C4 \texttt{realnewslike}
documents truncated to $300$ characters, which is dgMARK's published protocol; every method
in this paper uses the same construction, the same checkpoint and $256$ generated tokens.
Text quality is GPT-2-large perplexity, reported as a ratio against a non-watermarked
control \emph{produced under the same decoding regime}, since the three methods require
different regimes and a single shared baseline would charge the cost of a regime to
whichever watermark used it. Detectability is reported as the true positive rate at a
controlled false positive rate, never as a raw $z$.

% Baseline comparison for ReTrace. Style follows the C4 paper (booktabs, resizebox,
% oursemph / failred / finalres colour macros, deltas in parentheses).
% Requires in the preamble:
%   \usepackage[table]{xcolor}\usepackage{booktabs}
%   \definecolor{oursemph}{RGB}{0,90,181}
%   \definecolor{failred}{RGB}{178,24,43}
%   \definecolor{finalres}{RGB}{250,222,236}
%   \definecolor{notmeasured}{RGB}{178,24,43}
%
% Every number below is measured locally on one TITAN RTX: GSAI-ML/LLaDA-8B-Instruct in
% fp16, C4 realnewslike prompts truncated to 300 characters (dgMARK's published protocol),
% 256 generated tokens, perplexity under GPT-2-large, sequences of at least 200 tokens
% retained. TPR is reported at analytic FPR thresholds; the non-watermarked arm's measured
% FPR at the same thresholds is given in the notes so the thresholds can be audited.

\begin{table}[!t]
\centering
\small
\renewcommand{\arraystretch}{0.72}
\resizebox{\textwidth}{!}{%
\begin{tabular}{llrrrrrrr}
\toprule
& & & & \multicolumn{3}{c}{TPR at FPR $\uparrow$} & & \\
\cmidrule(lr){5-7}
Method & Setting & PPL $\downarrow$ & Ratio $\downarrow$ & $5\%$ & $1\%$ & $0.1\%$ & Gen.\ eval. & Det.\ eval. \\
\midrule
\multicolumn{9}{l}{\textit{Budget: one sequence evaluation per generated token}} \\
\midrule
dgMARK \citep{hong2026dgmark} & $k=1$, 256 tok & 15.40 & 1.55$\times$ & 0.88 & 0.74 & 0.62 & 256 & \textbf{0} \\
\rowcolor{finalres}
KGW \citep{kirchenbauer2023watermark} & $\delta=1$, 256 tok & 11.96 & \textbf{1.03}$\times$ & 0.97 & \textbf{0.93} & 0.73 & 256 & \textbf{0} \\
KGW \citep{kirchenbauer2023watermark} & $\delta=3$, 256 tok & 14.05 & 1.21$\times$ & 1.00 & 1.00 & 1.00 & 256 & \textbf{0} \\
Post-hoc substitution & best under $1.35\times$ PPL & 30.5 & 1.17$\times$ & 0.10 & 0.10 & 0.00 & 256 & 9 \\
\midrule
\multicolumn{9}{l}{\textit{Budget: $1.28\times$ decode --- ReTrace Response Bank plus reference schedule}} \\
\midrule
\textcolor{oursemph}{ReTrace} & $R{=}1$, 512 tok & 11.0 & \textbf{0.97}$\times$ & 0.57 & 0.43 & 0.20 & 656 & 144 \\
\textcolor{oursemph}{ReTrace} & $R{=}2$, 512 tok, held-out & -- & 1.14$\times$ & 0.83 & 0.73 & 0.63 & 656 & 144 \\
\textcolor{oursemph}{ReTrace} & $R{=}4$, 512 tok, held-out & -- & \textcolor{failred}{1.40$\times$} & 0.90 & 0.83 & 0.77 & 656 & 144 \\
\rowcolor{finalres}
\textcolor{oursemph}{ReTrace} & $R{=}16$, $\kappa{=}0.05$, 512 tok & -- & \textbf{0.74}$\times$ & 0.90 & \textbf{0.80} & 0.60 & 656 & 144 \\
\textcolor{oursemph}{ReTrace} & $R{=}8$, $\kappa{=}0.1$, 512 tok & 5.67 & \textbf{0.87}$\times$ & 0.80 & 0.70 & 0.60 & 656 & 144 \\
\rowcolor{finalres}
\textcolor{oursemph}{ReTrace} & $R{=}8$, $\kappa{=}0.1$, 1024 tok & -- & \textbf{0.83}$\times$ & 0.81 & \textbf{0.75} & \textbf{0.69} & 1312 & 288 \\
\midrule
\multicolumn{9}{l}{\textit{Budget: $\sim\!4\times$ decode (one-step lookahead per candidate)}} \\
\midrule
dgMARK \citep{hong2026dgmark} & $+3$-beam, 256 tok & 12.23 & 1.23$\times$ & \textbf{0.92} & \textbf{0.86} & \textbf{0.82} & \textcolor{failred}{$\sim$1024} & \textbf{0} \\
\end{tabular}%
}
\caption{\textbf{Detection, matched-control perplexity, and compute.} Rows are grouped by
generation budget and use one machine, checkpoint, and prompt protocol. \emph{Ratio} is PPL
relative to a non-watermarked control under the same decoding regime; a ratio below one
does not imply better semantic quality. \emph{Eval.} counts sequence evaluations, not
batched wall-clock invocations. ReTrace uses 144 extra evaluations to construct Response
Banks during a 512-token decode and 144 more for verification. dgMARK three-beam uses
approximately four generation evaluations per token. KGW is scored on distinct bigrams;
without deduplication, its unwatermarked control has $37\%$ FPR. dgMARK arms use $n=50$
and KGW arms $n=30$; principal 512-token ReTrace rows use $n=30$ ($n=20$ for $R=16$),
and the 1024-token row uses $n=16$. PPL ratios use only pairs passing the decoded-length rule.}
\label{tab:baselines}
\end{table}


Two observations about the baselines. Locally reproduced dgMARK is weaker than its
published table ($1.25\times$ for $99.4\,\%$), which is why the reproduction was worth
running rather than citing. And KGW's apparent $1.03\times$ is against its own
left-to-right control; charged against block decoding it is $1.20\times$, and under forced
left-to-right decoding this model repeats about a third of its bigrams, so its detector
must score each distinct bigram once --- an un-deduplicated detector measured a $37\,\%$
false positive rate on non-watermarked text.

% Statistical validity of the Shibboleth detector under the decision rule it actually ships.
\begin{table}[!t]
\centering
\small
\renewcommand{\arraystretch}{0.8}
\begin{tabular}{lrrl}
\toprule
Nominal $\alpha$ & Mean FPR & Maximum-key FPR & Relation \\
\midrule
$0.10$ & $0.0121$ & $0.0583$ & $\leq\alpha$ \\
$0.05$ & $0.0037$ & $0.0250$ & $\leq\alpha$ \\
$0.01$ & $0.0000$ & $0.0000$ & $\leq\alpha$ \\
\bottomrule
\end{tabular}
\caption{\textbf{Observed false-positive rates across keys.} $20$ keys $\times$
$120$ non-watermarked generations under the implemented decision rule.}
\label{tab:validity}
\end{table}


\paragraph{Held-out evaluation.}
The documents used while developing and debugging the method are not a test set. Final
numbers are reported on fresh documents with per-document nonces
$K_d = \mathrm{HMAC}(K, \mathrm{nonce}_d)$, so each document is an independent key draw and
the guarantee of Proposition~\ref{prop:null} applies per document.
